\documentclass[10pt]{article}

\usepackage{microtype,amsmath,amssymb,amsthm,booktabs}
\usepackage[margin=0.85in]{geometry}
\usepackage[colorlinks=true,linkcolor=blue,urlcolor=blue,citecolor=blue]{hyperref}
\usepackage{bookmark}
\newtheorem{theorem}{Theorem}
\newtheorem{lemma}[theorem]{Lemma}
\newtheorem{corollary}[theorem]{Corollary}
\newcommand{\dd}{\mathrm d}
\newcommand{\artanh}{\operatorname{artanh}}
\newcommand{\Tr}{\operatorname{Tr}}
\setlength{\emergencystretch}{2em}
\title{Reciprocal Traces and the Half-Filled Plunge Law}
\author{Kirtus Justus (Theoretical Physicist)\\
\small Port Coquitlam, British Columbia, Canada}
\date{October 2026}
\hypersetup{pdftitle={Reciprocal Traces and the Half-Filled Plunge Law},pdfauthor={Kirtus Justus (Theoretical Physicist)}}
\usepackage{fontspec,unicode-math,mathtools}
\setmainfont{DejaVu Serif}\setsansfont{DejaVu Sans}\setmonofont{DejaVu Sans Mono}\setmathfont{Latin Modern Math}
\begin{document}
\maketitle
\begin{abstract}
For the half-filled finite sine kernel at even size, the first reciprocal
trace of the deflated differential spectrum is an exact positive Laplace
transform. A weighted Legendre resolvent and a Mellin contour identity
determine its normalization. Taylor estimates on a fixed analytic kernel
then prove its large-size expansion to every fixed order, including an
explicit remainder after the term of order $N^{-9}$. This result yields
the signed plunge law for every nonzero Jacobi eigenmode whose eigenvalue
tends to zero. An elementary trial vector proves that the central pair
has this property. No uniform bulk quantum-period expansion is assumed.
\end{abstract}

\section{The finite problem and its exact differential trace}
Let $N\ge2$ be even. Define real symmetric matrices, with indices
$0\le j,k<N$, by
\[
B_{jk}=\begin{cases}1/2,&j=k,\\
\sin[\pi(j-k)/2]/[\pi(j-k)],&j\ne k,\end{cases}\qquad
H_{j-1,j}=H_{j,j-1}=\frac{j(N-j)}2,\quad H_{jj}=0,
\qquad J_N=\frac{2H}{N}.
\]
These are the usual half-filled concentration matrix and its commuting
Jacobi matrix \cite{eisler}. They commute, $0<B<I$, and their common
eigenpairs are denoted
\[
Hv=\beta v,\qquad J_Nv=\lambda v,\qquad Bv=\zeta v,
\qquad \beta=N\lambda/2.
\]
Put $\alpha=(N^2-1)/4$. On $(-\pi/2,\pi/2)$ consider the Friedrichs
realization
\[
A_Nu=-(\cos x\,u')'-\alpha\cos x\,u.
\]
Its eigenfunctions are regular, with no logarithmic term, at both
endpoints. Write its increasing eigenvalues as $\mu_{N,k}$, starting
at $k=0$, and define
\begin{equation}
Z_N=\sum_{k=N}^{\infty}\mu_{N,k}^{-1}.
\label{eq:Z}
\end{equation}
The exact spectral-tail construction \cite{tails} gives
\begin{equation}
\{\mu_{N,0},\ldots,\mu_{N,N-1}\}=-\operatorname{spec}H,
\qquad \mu_{N,N}>\|H\|,
\label{eq:finiteblock}
\end{equation}
and the norm-convergent scalar representation at every common mode
\begin{equation}
\log\frac\zeta{1-\zeta}=F_N(\beta),\qquad
F_N(b)=2\sum_{k=N}^{\infty}\artanh\frac b{\mu_{N,k}},
\quad |b|<\mu_{N,N}.
\label{eq:logit}
\end{equation}
For clarity, the ingredients of these exact inputs are as follows.
The centered finite Fourier space is invariant under $-A_N$ and supplies
$-\operatorname{spec}H$. A sign-change test against that space, together
with the variational nodal bound, identifies exactly the first $N$
differential eigenvalues. For the left regular solution normalized to
one, let $C_N(b)$ be its logarithmic coefficient at the right endpoint.
Its Green kernel gives
$C_N'/C_N=\Tr(A_N+b)^{-1}$. Deflating
$D_N(b)=\det(bI-H)$ removes precisely those first $N$ terms.
The endpoint Lagrange identity and the complementary Fourier arc give
the occupation ratio as $[C_N/D_N](\beta)/[C_N/D_N](-\beta)$.
Integrating the deflated resolvent trace proves \eqref{eq:logit}.
These statements concern the specified endpoint realization, not an
arbitrary self-adjoint extension.

At $b=0$ set $\nu=(N-1)/2$ and translate $s=x+\pi/2$.
The regular solutions are $P_\nu(\cos s)$ and $P_\nu(-\cos s)$,
where $P_\nu$ denotes the Ferrers function on $(-1,1)$.
The Legendre connection formula gives
$C_N(0)=2\sin(\pi\nu)/\pi$, while the parameter Lagrange identity gives
$C_N'(0)=-\int_0^\pi P_\nu(\cos s)P_\nu(-\cos s)\dd s$.
Because $D_N$ is even and nonzero at zero,
\begin{equation}
Z_N=\frac{C_N'(0)}{C_N(0)}
=-\frac\pi{2\sin\pi\nu}\int_0^\pi
P_\nu(\cos s)P_\nu(-\cos s)\dd s.
\label{eq:Legendre}
\end{equation}
The notation $P_{-1/2}(\cosh t)$ below instead denotes the real
Legendre function for argument at least one. All spectral and modular
quantities here are dimensionless; no physical clock or energy scale is
introduced.

\section{A weighted resolvent and a Mellin reflection identity}
Define
\[
\Phi(z)={}_2F_1(1/2,1/2;1;z),\qquad G(z)=\Phi(z)^2,
\qquad c_j=\frac{(1/2)_j}{j!},\qquad
a_l=\sum_{j=0}^l c_j^2c_{l-j}^2.
\]
Thus $G(z)=\sum_{l\ge0}a_lz^l$ for $|z|<1$.
\begin{lemma}[Weighted Legendre trace]
Let
\[
M(a)=\sum_{l=0}^{\infty}\frac{a_l}{l+a}.
\]
This defines a meromorphic function with possible poles only at the
nonpositive integers. For even $N$, with $\nu=(N-1)/2$,
\begin{equation}
Z_N=\frac\pi2\{M(-\nu)+M(\nu+1)\}.
\label{eq:Msum}
\end{equation}
For $\Re a>0$, $M(a)=\int_0^1z^{a-1}G(z)\dd z$.
\end{lemma}
\begin{proof}
The generating function for Legendre polynomials yields
\[
P_l(\cos s)=\sum_{j=0}^l c_jc_{l-j}e^{i(2j-l)s},
\qquad \int_0^\pi P_l(\cos s)^2\dd s=\pi a_l.
\]
Induction gives $c_j^2\le1/(j+1)$, and hence
\begin{equation}
0<a_l\le\frac{2H_{l+1}}{l+2},\qquad
H_r=\sum_{j=1}^r\frac1j.
\label{eq:acoef}
\end{equation}
Consider the ordinary Legendre operator with regular/Friedrichs endpoint
conditions, $L=-[(1-x^2)u']'$ on $L^2(-1,1;\dd x)$, with eigenvalues $l(l+1)$
and normalized eigenfunctions $\sqrt{(2l+1)/2}P_l(x)$.
Its shifted Green kernel has diagonal
\[
(L-\nu(\nu+1))^{-1}(x,x)
=-\frac\pi{2\sin\pi\nu}P_\nu(x)P_\nu(-x).
\]
Initially $-1/2<\nu<0$, so the resolvent is positive.
Integrating its diagonal with weight $(1-x^2)^{-1/2}$ gives, by
positive-kernel trace on compact intervals and monotone exhaustion,
\[
-\frac\pi{2\sin\pi\nu}\int_0^\pi
P_\nu(\cos s)P_\nu(-\cos s)\dd s
=\frac\pi2\sum_{l\ge0}a_l
\left(\frac1{l-\nu}+\frac1{l+\nu+1}\right).
\]
The series is locally normally convergent away from its poles by
\eqref{eq:acoef}. Endpoint logarithmic bounds, locally uniform in
$\nu$, give the meromorphic continuation of the integral expression.
Continuation to the required half-integer $\nu$ and
\eqref{eq:Legendre} prove \eqref{eq:Msum}. The Mellin integral for
$M$ follows by absolute convergence. At negative half-integers it
always means meromorphic continuation, not a divergent ordinary integral.
\end{proof}

\begin{lemma}[Mellin reflection at half-integers]
For $n\ge1$,
\begin{equation}
M(1/2-n)+M(1/2+n)
=\int_0^1 z^{n-1/2}G(1-z)\dd z.
\label{eq:reflection}
\end{equation}
\end{lemma}
\begin{proof}
For $0<a<1$, introduce
\[
J(c)=\int_0^1z^{c-1}G(1-z)\dd z,\qquad
Q(a)=\int_0^\infty x^{a-1}G(-x)\dd x.
\]
Use the principal branch of $\Phi$, cut along $[1,\infty)$.
Reciprocal-modulus continuation \cite{dlmf} gives, for $x>1$,
\[
\Phi(x+i0)=x^{-1/2}\{\Phi(1/x)+i\Phi(1-1/x)\},
\quad
\Re G(x+i0)=x^{-1}\{G(1/x)-G(1-1/x)\}.
\]
Integrate $z^{a-1}G(z)$ around the upper half-plane, with small
indentations at zero and one. The origin arc is $O(\epsilon^a)$;
the arc at one is $O(\epsilon(1+|\log\epsilon|)^2)$.
Since $\Phi(z)=O(|z|^{-1/2}(1+\log|z|))$ there, the large arc
is $O(R^{a-1}(1+\log R)^2)$ and also vanishes.
The negative-axis phase therefore gives
\[
\int_0^\infty x^{a-1}G(x+i0)\dd x=e^{i\pi a}Q(a).
\]
Taking real parts and changing $z=1/x$ on $x>1$ yields
\begin{equation}
M(a)+M(1-a)-J(1-a)=\cos(\pi a)Q(a).
\label{eq:contour}
\end{equation}
Pfaff's transformation gives
$G(-x)=(1+x)^{-1}G(x/(1+x))$, so
\[
Q(a)=\int_0^1z^{a-1}(1-z)^{-a}G(z)\dd z.
\]
For $\Re a<1$, its endpoint at one is integrable. At zero,
subtracting finitely many Taylor terms of $(1-z)^{-a}G(z)$
continues this integral meromorphically; its possible poles in that
half-plane are only $0,-1,-2,\ldots$.
Both sides of \eqref{eq:contour} consequently continue to $\Re a<1$.
At $a=1/2-n$, $Q(a)$ is finite and $\cos(\pi a)=0$.
This proves \eqref{eq:reflection}, including its normalization.
\end{proof}

\section{Exact positive representation and trace asymptotics}
\begin{theorem}[Positive Laplace formula]
For every positive even integer $N$,
\begin{equation}
\boxed{\quad Z_N=\pi\int_0^\infty e^{-Nt}f(t)\dd t,
\qquad f(t)=P_{-1/2}(\cosh t)^2.\quad}
\label{eq:laplace}
\end{equation}
Equivalently,
$Z_N=(\pi/2)\int_0^1 z^{(N-1)/2}G(1-z)\dd z$.
\end{theorem}
\begin{proof}
The equivalent form follows immediately from \eqref{eq:Msum} and
\eqref{eq:reflection}. The quadratic transformation
\[
y(t):=P_{-1/2}(\cosh t)=e^{-t/2}\Phi(1-e^{-2t})
\]
then gives \eqref{eq:laplace} on substituting $z=e^{-2t}$.
The transformation can be checked directly: both sides are the unique
regular solution of
\begin{equation}
y''+\coth t\,y'+\tfrac14 y=0,\qquad y(0)=1,\quad y'(0)=0.
\label{eq:yode}
\end{equation}
Thus no numerical match or asymptotic normalization is required.
\end{proof}

Euler's integral gives the particularly useful real formula
\begin{equation}
y(t)=\frac2\pi\int_0^{\pi/2}
\frac{\dd\varphi}{\sqrt{1+\sinh^2(t/2)\sin^2\varphi}}.
\label{eq:Euler}
\end{equation}
It follows that $0<f(t)\le1$, with strict decrease for $t>0$.
The right side of \eqref{eq:laplace} defines a completely monotone
interpolation for real $N>0$, but its identification with the spectral
trace is asserted here only at even integers. In particular, odd $N$
has a zero finite mode and requires a different deflation at $b=0$.

\begin{corollary}[An exact slope bound]
For even $N\ge2$,
\begin{equation}
0<\pi-NZ_N<\frac\pi{4N^2}.
\label{eq:slopebound}
\end{equation}
Moreover $NZ_N/\pi$ increases strictly along the positive even
integers and tends to one.
\end{corollary}
\begin{proof}
Equation \eqref{eq:yode} implies
\[
-y'(t)=\frac1{4\sinh t}\int_0^t\sinh s\,y(s)\dd s
<\tfrac14\tanh(t/2)<t/8\quad(t>0).
\]
Thus $0<1-f(t)=(1-y(t))(1+y(t))<t^2/8$.
Integration proves \eqref{eq:slopebound}.
Finally $NZ_N/\pi=\int_0^\infty e^{-u}f(u/N)\dd u$;
strict monotonicity and dominated convergence give the remaining claims.
\end{proof}

\begin{theorem}[Asymptotic expansion with a remainder]
As positive even $N\to\infty$,
\begin{equation}
\boxed{Z_N=\frac\pi N-\frac\pi{4N^3}
+\frac{17\pi}{64N^5}-\frac{169\pi}{256N^7}
+\frac{50777\pi}{16384N^9}+O(N^{-11}).}
\label{eq:traceexp}
\end{equation}
There is an ordinary asymptotic expansion to every fixed order.
In \eqref{eq:traceexp}, the absolute remainder is at most
$\pi B\,10!/N^{11}$, where
\begin{align*}
\rho&=3/2,& M&=(1-\sinh^2(3/4))^{-1},\\
B&=\max\left\{\frac{M\rho^{-10}}{1-\rho^{-2}},\,
2+\frac18+\frac{17}{1536}+\frac{169}{184320}
+\frac{50777}{660602880}\right\}.
\end{align*}
\end{theorem}
\begin{proof}
Expansion of \eqref{eq:yode}, or of its defining hypergeometric series,
gives
\begin{equation}
f(t)=1-\frac{t^2}{8}+\frac{17t^4}{1536}
-\frac{169t^6}{184320}+\frac{50777t^8}{660602880}+O(t^{10}).
\label{eq:fseries}
\end{equation}
For complex $|t|\le\rho$, $|\sinh(t/2)|\le\sinh(3/4)<1$.
Equation \eqref{eq:Euler}, with the square root continued from zero,
shows that $f$ is even and analytic on a neighborhood of this disk,
with $|f|\le M$. Cauchy's estimate therefore bounds its coefficient
of $t^{2j}$ by $M\rho^{-2j}$.
Let $S_8$ be the polynomial displayed in \eqref{eq:fseries}.
For $0\le t\le1$, the geometric coefficient tail gives
$|f(t)-S_8(t)|\le M\rho^{-10}t^{10}/(1-\rho^{-2})$.
For $t\ge1$, $0<f\le1$ gives the other term defining $B$.
Thus $|f-S_8|\le Bt^{10}$ on the entire positive axis.
Integrating this finite Taylor estimate in \eqref{eq:laplace} proves
the stated remainder and coefficients. The same argument applies at
each fixed truncation order. It does not assert convergence of the
infinite inverse-power series.
\end{proof}

The coefficient of $N^{-9}$ is now a spectral-trace theorem.
Its agreement with a formal quantum-period coefficient does not,
by itself, prove that period expansion at other spectral modes.

\section{A uniform estimate near the center}
Define the leading reference function and errors by
\begin{equation}
\eta(t)=\pi t\,{}_3F_2(1/4,1/2,3/4;1,3/2;4t^2),\qquad
\varepsilon=\log\frac{1-\zeta}{\zeta},\qquad
e=\varepsilon+N\eta(\lambda/N).
\label{eq:errors}
\end{equation}
\begin{theorem}[Signed plunge law for shrinking Jacobi modes]
Let even $N\to\infty$ and choose nonzero common eigenmodes with
$\lambda_N\to0$. Then
\begin{equation}
e_N=\frac{\pi\lambda_N}{4N^2}
\left[1+O(N^{-2}+\lambda_N^2)\right],\qquad
\varepsilon_N=-\pi\lambda_N\left[1+O(N^{-2})\right].
\label{eq:plunge}
\end{equation}
Consequently
\begin{equation}
\boxed{\quad\frac{-4N^2e_N}{\varepsilon_N}\longrightarrow1.\quad}
\label{eq:ratio}
\end{equation}
The estimates apply to either sign of the eigenvalue.
\end{theorem}
\begin{proof}
We give a bound uniform for even $N\ge8$ and $0<|\lambda|\le1$.
The central two-site Rayleigh quotient gives $\|H\|\ge N^2/8$.
By \eqref{eq:finiteblock},
\[
r:=\frac{|\beta|}{\mu_{N,N}}\le\frac{4|\lambda|}{N}\le\frac12,
\qquad
\sum_{k\ge N}\mu_{N,k}^{-3}
\le\frac{64Z_N}{N^4}\le\frac{64\pi}{N^5}.
\]
The positive artanh series in \eqref{eq:logit} gives
\begin{equation}
|F_N(\beta)-N\lambda Z_N|
\le\frac{2|\beta|^3}{3(1-r^2)}\sum_{k\ge N}\mu_{N,k}^{-3}
\le\frac{64\pi|\lambda|^3}{9N^2}.
\label{eq:nonlinear}
\end{equation}
The coefficients $d_k$ of the hypergeometric series in \eqref{eq:errors}
are positive, with $d_1=1/16$ and $d_{k+1}<d_k$.
Hence, for $|t|<1/2$,
\[
|\eta(t)-\pi t|\le\frac{\pi|t|^3}{4(1-4t^2)},
\qquad
|N\eta(\lambda/N)-\pi\lambda|
\le\frac{4\pi|\lambda|^3}{15N^2}.
\]
Since $\varepsilon=-F_N(\beta)$, these inequalities imply
\begin{equation}
|e-\lambda(\pi-NZ_N)|\le\frac{8\pi|\lambda|^3}{N^2}.
\label{eq:uniform}
\end{equation}
Equation \eqref{eq:traceexp} gives
$\pi-NZ_N=\pi/(4N^2)+O(N^{-4})$.
This and \eqref{eq:uniform} prove the first estimate in
\eqref{eq:plunge}; \eqref{eq:slopebound} and
\eqref{eq:nonlinear} prove the second. Dividing is legitimate:
\eqref{eq:logit} is strictly increasing and vanishes only at $\beta=0$.
No differentiation of an unproved spectral asymptotic has been used.
\end{proof}

\section{The actual central pair satisfies the hypothesis}
\begin{lemma}[An elementary central-mode bound]
Let $N=2m$, $m\ge1$, and let $\lambda_{\min,+}(J_{2m})$ be the
smallest positive eigenvalue. With $H_0=0$,
\begin{equation}
0<\lambda_{\min,+}(J_{2m})
\le\frac{2m-1}{2m}\left(1+\frac14H_{m-1}\right)^{-1/2}.
\label{eq:centralbound}
\end{equation}
In particular the innermost positive and negative modes satisfy
\eqref{eq:plunge} and \eqref{eq:ratio}.
\end{lemma}
\begin{proof}
Use site indices $1,\ldots,2m$ and bonds
$b_j=j(2m-j)/(2m)$. Set $v_{2k}=0$, $v_{2k-1}=A_k$,
$A_1=1$, and $A_{k+1}=-b_{2k-1}A_k/b_{2k}$.
All rows of $J_{2m}v$ vanish except the last, so
$\|J_{2m}v\|=b_{2m-1}|A_m|$.
With $c_r=(1/2)_r/r!$ as above,
\[
|A_k|=\prod_{j=1}^{k-1}\frac{2j-1}{2j}
\frac{2m-2j+1}{2m-2j}\ge c_{k-1},\qquad
|A_m|=(2m-1)c_{m-1}^2\le1.
\]
For the last inequality, the ratio of successive endpoint amplitudes
is $1-1/(4m^2)<1$, starting at one.
Also $r c_r^2$ increases for $r\ge1$, starting at $1/4$,
because its successive ratio is $(2r+1)^2/[4r(r+1)]>1$.
Thus $\|v\|^2\ge1+H_{m-1}/4$.
Expanding $v$ in an orthonormal eigenbasis gives
$\operatorname{dist}(0,\operatorname{spec}J_{2m})
\le\|J_{2m}v\|/\|v\|$.
Finally the spectrum is symmetric, and
$\det J_{2m}=(-1)^m\prod_{k=1}^m b_{2k-1}^2\ne0$.
The distance is therefore the smallest positive eigenvalue, proving
\eqref{eq:centralbound}. The harmonic sum diverges, so the required
nonzero eigenvalues tend to zero.
\end{proof}
The estimate is sufficient but is not claimed to give the sharp
central eigenvalue scale. In combination with the previous theorem
it gives the conservative central-pair bound
$-4N^2e_N/\varepsilon_N=1+O(1/\log N)$.

\section{Independent recurrence and scope}
An additional exact characterization, useful for checking constants,
is
\begin{equation}
(N+1)^3Z_{N+2}-2N(N^2+1)Z_N+(N-1)^3Z_{N-2}=0,
\qquad N\ge4\ \text{even}.
\label{eq:recurrence}
\end{equation}
For a direct derivation from \eqref{eq:Legendre}, let
$I_\nu=\int_{-1}^1(1-x^2)^{-1/2}P_\nu(x)P_\nu(-x)\dd x$
and insert an extra $x^2$ to define $T_\nu$.
The degree raising and lowering identities \cite{dlmf}, followed
by two integrations by parts, give
\begin{align*}
I_{\nu+1}&=\frac\nu{(\nu+1)^2}T_\nu
+\left[-1+\frac1{2(\nu+1)^2}\right]I_\nu,\\
I_{\nu-1}&=-\frac{\nu+1}{\nu^2}T_\nu
+\left[-1+\frac1{2\nu^2}\right]I_\nu.
\end{align*}
Endpoint products are $O(r^{1/2}(1+|\log r|))$ and vanish.
Eliminating $T_\nu$ gives
\[
(\nu+1)^3I_{\nu+1}
+(2\nu+1)(\nu^2+\nu+1/2)I_\nu+\nu^3I_{\nu-1}=0.
\]
Substitution in \eqref{eq:Legendre}, including the alternating sine
factor, proves \eqref{eq:recurrence}. The recurrence alone would not
have fixed the normalization of a candidate solution; the contour
argument supplies that missing information.

Numerical evaluation of the connection derivative and
both positive-integral forms agrees at $N=2,4,8$ to 70-digit working
precision. Exact rational checks verify the displayed kernel
coefficients and trial-vector identities. These calculations support
reproducibility; the proofs and remainder estimates above do not
depend on finite numerical agreement.

The conclusions concern even size and half filling. They prove a
reciprocal-trace expansion and the signed law for shrinking Jacobi
modes, including the central pair. They do not establish the uniform
bulk quantum-period remainder, its higher spectral coefficients,
a general-filling effective-size convention, or the full Wilson
relative-entropy finite-part asymptotic. The separate companion \emph{Uniform Bulk Modular Asymptotics and the Leading Wilson Entropy Law} proves the fixed-interior period expansion and leading Wilson law with bounded remainder.

\paragraph{Acknowledgments and provenance.}
This work belongs to \emph{From Prephysical Hypotheses to Physical Models}.
The analysis and checks here were developed by Kirtus Justus with
ChatGPT (OpenAI). The preceding connection-coefficient proposal was
supplied in correspondence prepared with Claude (Anthropic).
The wider program has also used Gemini (Google), Grok (SpaceXAI),
and DeepSeek (High-Flyer); no specific proof here is attributed to
those systems. Assistance does not imply external peer review or
endorsement. No priority claim is made for the special-function identities.

\begin{thebibliography}{9}
\bibitem{tails} K. Justus, \emph{Spectral Tails and Exact Concentration Logits},
research manuscript, October 2026.
\bibitem{eisler} V. Eisler and I. Peschel, Analytical results for the
entanglement Hamiltonian of a free-fermion chain,
\emph{J. Phys. A} \textbf{50} (2017), 284003,
\href{https://arxiv.org/abs/1703.08126}{arXiv:1703.08126}.
\bibitem{dlmf} NIST Digital Library of Mathematical Functions,
\href{https://dlmf.nist.gov/14.3}{\S14.3},
\href{https://dlmf.nist.gov/14.10}{\S14.10},
\href{https://dlmf.nist.gov/15.6}{\S15.6},
\href{https://dlmf.nist.gov/15.8}{\S15.8}, and
\href{https://dlmf.nist.gov/19.7}{\S19.7}; in particular
(14.10.4)--(14.10.5), (15.6.1), (15.8.1), and (19.7.3).
\end{thebibliography}
\end{document}
